% Lente convergente con f = 12 cm e oggetto a p = 60 cm: immagine reale a q = 15 cm,
% fra F' e 2F'. Scala: 0,12 cm di figura per 1 cm reale; altezza dell'oggetto
% arbitraria, immagine alta un quarto (m = -1/4).
% Esercizio ud7-ottica-coniugati (soluzione)
\begin{tikzpicture}[>=Stealth, font=\small]
  \draw[->] (-7.9,0) -- (3.4,0) node[right] {asse};
  \draw[<->, very thick, formula] (0,-1.4) -- (0,1.9) node[above] {lente};
  \foreach \x/\nome in {-2.88/2F, -1.44/F} {
    \fill (\x,0) circle (1.3pt) node[below] {$\nome$};
  }
  \foreach \x/\nome in {1.44/F', 2.88/2F'} {
    \fill (\x,0) circle (1.3pt) node[above right=-1pt] {$\nome$};
  }
  % oggetto e immagine
  \draw[->, very thick] (-7.2,0) -- (-7.2,1.6) node[above] {oggetto};
  \draw[->, very thick, attenzione] (1.8,0) -- (1.8,-0.4);
  \node[attenzione, anchor=north east] at (1.75,-0.45) {immagine};
  % raggi notevoli
  \draw[esempio, thick] (-7.2,1.6) -- (0,1.6) -- (2.3,-0.956);
  \draw[medica, thick] (-7.2,1.6) -- (2.3,-0.511);
  \draw[approfondimento, thick] (-7.2,1.6) -- (0,-0.4) -- (2.3,-0.4);
  % distanze
  \draw[|<->|] (-7.2,-1.1) -- (0,-1.1) node[midway, below] {$p = \SI{60}{\centi\metre}$};
  \draw[|<->|] (0,-1.1) -- (1.8,-1.1) node[right=2pt] {$q = \SI{15}{\centi\metre}$};
\end{tikzpicture}
